%%%%%%%%%%%%%%%%%%%%%%%%%%%%%%%%%%%%%%%%%%%%%%%%%%%%%%%
\begin{frame}

\frametitle{Including Files}

The {\texttt{Snakefile}} can be organized better using the {\blue{include}} directive:
\vs

\begin{figure}
%   \centering
  \begin{tabular}{lc}
  include:  "accessors.py" \\
  \end{tabular}
\end{figure}

\end{frame}


%%%%%%%%%%%%%%%%%%%%%%%%%%%%%%%%%%%%%%%%%%%%%%%%%%%%%%%
\begin{frame}

\frametitle{Constraints}

If a file specification has a mix of variables (i.e., {\texttt{\{dataset\}-\{count\}}}, then Snakemake cannot parse the filename {\texttt{iris-2-42}}.  One solution is to employ constraints on variables:

\begin{figure}
%   \centering
  \begin{tabular}{lc}
  wildcard\_constraints: \\
  {\D\D}count = "{\textbackslash}d+" \\
  \end{tabular}
\end{figure}

Python regular expressions are allowed.  Thus, it might even be better to say {\texttt{dataset}} cannot have hyphens and use underscores instead.

\end{frame}


%%%%%%%%%%%%%%%%%%%%%%%%%%%%%%%%%%%%%%%%%%%%%%%%%%%%%%%
\begin{frame}

\frametitle{Configuration File}

Allows the specification of a YAML file, which is then read in as a global dictionary.  If {\texttt{config.yaml}} looks like this:
\vs

\begin{figure}
%   \centering
  \begin{tabular}{lc}
  datasets:\\
  {\D\D}iris:\\
  {\D\D}{\D\D}numcols: 5\\
  \end{tabular}
\end{figure}

Read it in like this in your {\texttt{Snakefile}}:

\begin{figure}
%   \centering
  \begin{tabular}{lc}
  configfile: "config.yaml" \\
  \end{tabular}
\end{figure}

And then access it like this:

\begin{figure}
%   \centering
  \begin{tabular}{lc}
  curr\_sample = "iris"\\
  numcols = config['datasets'][curr\_sample]['numcols']\\
  \end{tabular}
\end{figure}

\end{frame}


%%%%%%%%%%%%%%%%%%%%%%%%%%%%%%%%%%%%%%%%%%%%%%%%%%%%%%%
\begin{frame}

\frametitle{Suggestions}

\begin{itemize}
  \item Suggest each rule write to a {\green{different}} directory to avoid the possibility of a circular dependency.
  \item If a shell command in a rule returns a non-zero value, the rule fails and (for safety reasons) all outputs are deleted since they are inconsistent.
  \item Log files for programs should {\red{not}} be specified in the {\blue{output}} directive.
\end{itemize}

\end{frame}


%%%%%%%%%%%%%%%%%%%%%%%%%%%%%%%%%%%%%%%%%%%%%%%%%%%%%%%
\begin{frame}

\frametitle{Running Snakemake}

\begin{itemize}
  \item After running {\texttt{Snakemake}}, a ``hidden'' directory called {\texttt{.snakemake}} is created which keeps track of its progress.  If you need to reset things, then just erase this directory.
  \item If a previous run was interrupted, then use the {\texttt{--unlock}} argument to remove the lock on the working directory.
\end{itemize}

\end{frame}


%%%%%%%%%%%%%%%%%%%%%%%%%%%%%%%%%%%%%%%%%%%%%%%%%%%%%%%
\begin{frame}

\frametitle{Running Snakemake}

In fact, Snakemake has many useful command-line arguments:
    \begin{itemize}
      \item {\texttt{--configfile}}  Explicitly specify the configuration file.
      \item {\texttt{--cores}}  Set the number of cores to use.
      \item {\texttt{--printshellcmds}} / {\texttt{-p}}  Print shell commands as they are executed.
      \item {\texttt{--dag}}  Output a directed acyclic graph in the dot language.
      \item {\texttt{--dryrun}}  Perform a dry run.
      \item {\texttt{--list}}  List rules in the Snakefile.
      \item {\texttt{--forceall}} / {\texttt{-F}}  Force the execution of a rule and all subsequent rules.
    \end{itemize}

\end{frame}

